\documentclass[10pt]{article}
\usepackage[T1]{fontenc}
\usepackage[utf8]{inputenc}
\usepackage[french]{babel}
\usepackage{amsmath,amssymb,mathrsfs,mathpazo}
\usepackage{geometry}
\geometry{a4paper,margin=2cm}
\usepackage{graphicx,array,tabularx,tabularray,colortbl,multirow,multicol,booktabs,enumitem,pifont,ulem,fancybox,diagbox,cancel}
\usepackage{pstricks,pstricks-add,pst-plot,pst-tree,pst-node,pst-text,pst-func,pst-math,pst-bspline,pst-3dplot,pst-eucl,pst-solides3d}
\let\pscalculate\undefined
\usepackage{pst-bezier}
\usepackage{tikz,pgfplots,tkz-tab}
\usetikzlibrary{arrows,arrows.meta,patterns,calc,positioning,trees,decorations.markings,shapes.geometric,shapes.misc,shapes.arrows,intersections,decorations.pathmorphing,decorations.pathreplacing,backgrounds,shadows,babel}
\pgfplotsset{compat=1.18}
\usepgfplotslibrary{fillbetween}
\usepackage[np]{numprint}
\usepackage{listings,makecell,pas-tableur,twoopt}
\let\barre\undefined
\usepackage[pstricks]{tablvar}
\definecolor{darkgreen}{RGB}{0,100,0}
\providecommand{\red}{\color{red}}
\providecommand{\blue}{\color{blue}}
\providecommand{\green}{\color{green}}
\providecommand{\black}{\color{black}}
\providecommand{\white}{\color{white}}
\providecommand{\gray}{\color{gray}}
\providecommand{\cyan}{\color{cyan}}
\providecommand{\magenta}{\color{magenta}}
\providecommand{\yellow}{\color{yellow}}
\providecommand{\violet}{\color{violet}}
\providecommand{\orange}{\color{orange}}
\providecommand{\euro}{\texteuro}
\providecommand{\EUR}{\texteuro}
\providecommand{\ieme}{\textsuperscript{e}}
\providecommand{\degres}{\ensuremath{^\circ}}
\providecommand{\degree}{\ensuremath{^\circ}}
\providecommand{\vv}[1]{\overrightarrow{#1}}
\providecommand{\Coord}[1]{#1}
\providecommand{\decosix}{}
\providecommand{\starredbullet}{\ensuremath{\bullet}}
\pagestyle{empty}
\setlength{\parindent}{0pt}
\usepackage{adjustbox}\DeclareUnicodeCharacter{2236}{:}\DeclareUnicodeCharacter{2190}{\ensuremath{\leftarrow}}
\providecommand{\R}{\mathbb{R}}
\providecommand{\N}{\mathbb{N}}
\providecommand{\Z}{\mathbb{Z}}
\providecommand{\Q}{\mathbb{Q}}
\providecommand{\D}{\mathbb{D}}
\providecommand{\C}{\mathbb{C}}
\providecommand{\e}{{\mathrm{e}}}
\providecommand{\vect}[1]{\overrightarrow{#1}}
\providecommand{\vectt}[1]{\overrightarrow{\mathrm{#1}}}
\providecommand{\vvt}[1]{\overrightarrow{\mathrm{#1}}}
\providecommand{\np}[2][]{{#2~#1}}
\providecommand{\vv}[1]{{\overrightarrow{#1}}}
\renewcommand{\d}{{\mathrm{d}}}
\providecommand{\degres}{{^\circ}}
\providecommand{\degre}{{^\circ}}
\providecommand{\blue}{\color{blue}}
\providecommand{\red}{\color{red}}
\providecommand{\enskip}{\quad}
\providecommand{\llbracket}{\lbrack\!\lbrack}
\providecommand{\rrbracket}{\rbrack\!\rbrack}
\providecommand{\euro}{{\text{€}}}
\providecommand{\no}{{\text{n°}}}
\providecommand{\interval}[2]{{[#1;#2]}}
\providecommand{\intervFO}[2]{{[#1;#2[}}
\providecommand{\dotfill}{\ldots}

\providecommand{\ds}{\displaystyle}
\providecommand{\textding}[1]{#1}
\providecommand{\Oij}{\(\mathrm{O};\vec{\imath},\vec{\jmath}\)}
\providecommand{\Oijk}{\(\mathrm{O};\vec{\imath},\vec{\jmath},\vec{k}\)}
\providecommand{\Ouv}{\(\mathrm{O};\vec{u},\vec{v}\)}
\providecommand{\intervFF}[2]{\ensuremath{[#1;#2]}}
\renewcommand{\euro}{\ifmmode\text{\texteuro}\else\texteuro\fi}
\DeclareUnicodeCharacter{00B0}{\ifmmode^{\circ}\else\textdegree\fi}
\DeclareUnicodeCharacter{20AC}{\ifmmode\text{\texteuro}\else\texteuro\fi}
% Énoncés : Éducation nationale. Transcription et corrigés : APMEP, Denis Vergès et contributeurs. Découpage et classement : CoopMaths, Sébastien Lozano, Éric Elter et contributeurs. Adaptation de la présentation : Maths Explorer.
% CC BY-SA : https://coopmaths.fr/www/about/
% https://coopmaths.fr/alea/static/e3c/2020/tex/e3c_2020_00_sujet26_2_cor.tex
\begin{document}
\noindent\textit{APMEP / CoopMaths. CC BY-SA. Présentation adaptée par Maths Explorer.}\par\medskip
\[
	f(x) = 8x^3 - 6x^2 - 2.
	\]
	
	\subsection*{1.a}
	\[
	f(x) = (x - 1)\left(8x^2 + 2x + 2\right).
	\]
	Développons :
	\begin{align*}
		(x - 1)\left(8x^2 + 2x + 2\right) &= x^3 + 2x^2 + 2x - 8x^2 - 2x - 2 \\
		&= x^3 - 6x^2 - 2 = f(x).
	\end{align*}
	La factorisation est juste.
	
	\subsection*{1.b}
	Un point de l'axe des abscisses a une ordonnée nulle, donc un point commun à cet axe et à \(C\), vérifie :
	\begin{align*}
		&8x^3 - 6x^2 - 2 = 0 \\
		\iff \quad & (x - 1)\left(8x^2 + 2x + 2\right) = 0 \\
		\iff \quad & 
		\begin{cases}
			x - 1 = 0 \\
			8x^2 + 2x + 2 = 0
		\end{cases}
	\end{align*}
	La première équation a pour solution \(x = 1\), la seconde est une équation du second degré pour laquelle :
	\[
	\Delta = 36 - 4 \times 8 \times 2 = 36 - 64 = -28 < 0.
	\]
	Cette équation n'a pas de solution réelle.
	Donc, \(C\) coupe l'axe des abscisses au seul point \((1\,;\,0)\).
	
	\subsection*{2.a}
	\(f\), fonction polynôme, est dérivable sur \(\mathbb{R}\) et sur cet intervalle :
	\[
	f'(x) = 24x^2 - 12x = 12x(2x - 1).
	\]
	
	\subsection*{2.b}
	Le trinôme \(x(2x - 1)\) a deux racines : \(0\) et \(\dfrac{1}{2}\).
	
	Le coefficient \(a = 24 > 0\), donc le trinôme est positif sauf sur l'intervalle \(\left]0\,;\,\dfrac{1}{2}\right[\).
	
	On en déduit le tableau de variations :
	\begin{center}
	\begin{tikzpicture}
		\tkzTabInit[lgt=2.5, espcl=2]{$x$ / 1, {Signe de $f'(x)$} / 1, {$f$} / 2}{${-\infty}$, ${0}$, ${\frac12}$, ${+\infty}$}
		\tkzTabLine{,+,0,-,0,+,}
		\tkzTabVar{-/{$ $},+/{$-2$},-/{$-\frac52$},+/{$ $}}{/}
	\end{tikzpicture}
\end{center}
		\(f(0) = -2\) est un maximum local de \(f\), et :
		\[
		f\left(\dfrac{1}{2}\right) = 8\left(\dfrac{1}{2}\right)^3 - 6\left(\dfrac{1}{2}\right)^2 - 2 = 1 - \dfrac{3}{2} - 2 = -\dfrac{5}{2},
		\]
		est un minimum local de la fonction \(f\).

	
	\subsection*{3.}
	Une équation de la tangente \( \mathcal T \) est :
	\[
	M(x\,;\,y) \in  \mathcal T \iff y - f\left(\dfrac{1}{2}\right) = f'\left(\dfrac{1}{2}\right)\left(x - \dfrac{1}{2}\right).
	\]
	Avec \( f\left(\dfrac{1}{2}\right) = -\dfrac{5}{2} \) et \( f'\left(\dfrac{1}{2}\right) = 12 \times \dfrac{1}{2} \left(2 \times \dfrac{1}{2} - 1\right) = 6 \times 0 = 0 \), l'équation devient :
	\[
	M(x\,;\,y) \in  \mathcal T \iff y + \dfrac{5}{2} = 0 \left(x - \dfrac{1}{2}\right) \iff y = -\dfrac{5}{2}.
	\]
	Comme $ B\left(0\,;\,-\dfrac{5}{2}\right) \in  \mathcal T$, on peut dire que le point \( B \) appartient à \(  \mathcal T \).

\end{document}
